% ============================================================================
% Sous-section à placer à la fin de la section 8 (Résultats réels), après
% l'investigation sur les e2-medium. Chiffres : results/exp_newvms/report.md
% (scripts/exp_newvms.sh puis scripts/newvms_report.sh).
% ============================================================================
\subsection{Généralisation : des VM jamais mesurées}
\label{sec:generalisation}

Notre modèle est calibré sur les quatre VM qui servent aussi aux expériences. Il
ne prouve donc pas qu'il sait prédire \emph{un autre} parc de machines. Pour le
vérifier, nous avons créé quatre VM jamais mesurées, de types et de familles
différents, et prédit leurs temps \textbf{avant} de regarder leurs mesures :

\begin{itemize}
  \item \textbf{e2-standard-4} : même famille que notre parc, mais 4~vCPU ;
  \item \textbf{n2-standard-2} : autre famille, processeur Intel à 2,8~GHz ;
  \item \textbf{n2d-standard-2} : autre famille, processeur AMD EPYC 7B13 ;
  \item \textbf{t2d-standard-2} : AMD EPYC 7B13, sans \emph{hyperthreading}
        (1~vCPU = 1~cœur physique).
\end{itemize}

Round Robin a exécuté les lots de 500 et 2000 tâches sur ces quatre VM, 3 fois
chacun (\mesure{}, \code{results/exp\_newvms/}). Les mesures sont très stables
($\pm 0{,}3\,\%$), sans bridage. Nous avons comparé la fin mesurée de chaque VM à
cinq prédictions, de la moins coûteuse à la plus coûteuse en mesures :

\begin{itemize}
  \item \textbf{A, fiche seule} : benchmark de vm2 uniquement, la nouvelle VM
        n'est décrite que par sa fiche (nombre de vCPU, RAM, prix) ;
  \item \textbf{B, fiche + 1 tâche} : une seule tâche mesurée sur la nouvelle VM
        donne un facteur de vitesse unique ;
  \item \textbf{D, fiche + 4 tâches} : une tâche par type de workload (taille
        moyenne), soit un facteur de vitesse par type, environ 2~s de mesure ;
  \item \textbf{E, comme D, en tenant compte de l'hyperthreading} : contention
        de 1,00 au lieu de 1,40 pour une VM dont chaque vCPU est un cœur entier ;
  \item \textbf{C, benchmark complet} : les 12 combinaisons mesurées sur la VM,
        comme pour notre parc (environ 3~min).
\end{itemize}

\begin{table}[H]
\centering\small
\begin{tabular}{llccccc}
\toprule
\textbf{VM} & \textbf{Processeur} & \textbf{A} & \textbf{B} & \textbf{D} & \textbf{E} & \textbf{C} \\
\midrule
e2-standard-4  & Intel 2,2~GHz  & \textbf{8\,\%}  & 6\,\%  & 4\,\%  & 4\,\%  & 5\,\%  \\
n2-standard-2  & Intel 2,8~GHz  & 27\,\% & 10\,\% & 8\,\%  & 8\,\%  & 9\,\%  \\
n2d-standard-2 & AMD EPYC 7B13  & 40\,\% & 83\,\% & 10\,\% & 10\,\% & 10\,\% \\
t2d-standard-2 & AMD EPYC 7B13  & 55\,\% & 39\,\% & 33\,\% & \textbf{14\,\%} & 34\,\% \\
\bottomrule
\end{tabular}
\caption{Erreur moyenne de la fin prédite de chaque VM jamais mesurée, sur
6 exécutions réelles (\mesure{}). A : fiche seule ; B : + 1 tâche ; D : + 4 tâches ;
E : + 4 tâches et hyperthreading ; C : benchmark complet.}
\label{tab:generalisation}
\end{table}

\begin{table}[H]
\centering\small
\begin{tabular}{lcccc}
\toprule
\textbf{VM} & \textbf{cpu\_task} & \textbf{matrix\_task} & \textbf{memory\_task} & \textbf{sort\_task} \\
\midrule
e2-standard-4  & 1,03 & 1,07 & 1,07 & 1,03 \\
n2-standard-2  & 1,29 & 1,30 & 1,35 & 1,30 \\
n2d-standard-2 & \textbf{4,66} & 1,69 & 1,57 & 1,43 \\
t2d-standard-2 & \textbf{4,69} & 1,76 & 1,59 & 1,43 \\
\bottomrule
\end{tabular}
\caption{Vitesse d'une tâche seule, relative à vm2 (taille moyenne de chaque
type, \mesure{}). Sur une même VM, ce facteur ne dépend presque pas de la taille.}
\label{tab:vitesse_type}
\end{table}

\paragraph{Lecture.}
\begin{itemize}
  \item \textbf{Dans la même famille, la fiche suffit} : l'e2-standard-4, jamais
        mesurée et avec deux fois plus de vCPU, est prédite à 8\,\% près à partir du
        seul benchmark de vm2. Son ralentissement réel en parallèle (1,42) est
        celui que le modèle déduit de son nombre de vCPU (1,40). Le modèle ne colle
        donc pas aux instances mesurées : il dépend du \emph{type} de machine.
  \item \textbf{Hors de la famille, la fiche ne suffit pas} : la vitesse du
        processeur n'y figure pas. Les nouvelles VM étant plus rapides que vm2, la
        prédiction A est trop pessimiste (27 à 55\,\%).
  \item \textbf{Une seule tâche de calibration peut induire en erreur.} Sur
        processeur AMD, \code{cpu\_task} va 4,7 fois plus vite que sur vm2, mais
        les autres types seulement 1,4 à 1,8 fois. Calibrer sur \code{cpu\_task}
        seule surestime la vitesse (B : 83\,\% d'erreur sur n2d).
  \item \textbf{Quatre tâches suffisent} : une par type de workload (D) donne la
        même précision que le benchmark complet (C), 4 à 10\,\% sur trois VM sur
        quatre, pour environ 2~s de mesure au lieu de 3~min.
  \item \textbf{L'hyperthreading doit faire partie de la fiche} : sur t2d, chaque
        vCPU est un cœur entier, et deux tâches simultanées ne se ralentissent pas
        (contention réelle 0,90). En le prenant en compte (E), l'erreur sur t2d
        passe de 33 à 14\,\%.
\end{itemize}

\paragraph{Conclusion.} Le modèle se généralise à un parc jamais mesuré à
condition de connaître, pour chaque type de machine, son nombre de vCPU, la
présence d'hyperthreading et une vitesse par type de tâche, que quatre tâches de
calibration suffisent à estimer. Un scheduler en production pourrait lancer ces
quatre tâches au démarrage de chaque nouvelle VM, en quelques secondes, au lieu
d'un benchmark complet.
